\documentclass[11pt]{article}
\usepackage[a4paper,margin=25mm]{geometry}
\usepackage{amsmath,amssymb}
\usepackage[T1]{fontenc}
\usepackage{lmodern}
\usepackage{hyperref}
\setlength{\parindent}{0pt}
\setlength{\parskip}{7pt}
\title{A prime-order Latin square with no transversal}
\author{SucRunBug}
\date{2 October 2026}
\begin{document}
\maketitle
\textbf{Conjecture 00000000058 -- FALSE.}

\textbf{Filed claim addressed.} Every Latin square of prime order $p$ has at least $c^p$ transversals for an absolute constant $c>1$.

\section*{Proof}
Use the Latin square of order $2$ given by addition modulo $2$:
\[\begin{pmatrix}0&1\\1&0\end{pmatrix}.\]
Its rows and columns are permutations of the two symbols. A transversal must choose one cell from each row and each column, with distinct symbols. There are exactly two column permutations: the diagonal chooses symbols $0,0$, and the off-diagonal chooses $1,1$. Neither is a transversal. The number of transversals is therefore zero. The order $p=2$ is prime.

For every proposed $c>1$, $c^2>0$, so the asserted lower bound fails. The statement does not exclude the even prime or require sufficiently large primes. This proof does not settle a version restricted to odd or sufficiently large primes.

\section*{Formalization scope}
Lean formalizes Latin rows and columns as bijections, transversals as column permutations with injective symbol lists, and computes zero transversals for the order-two square. It proves failure for every c>1.

\textbf{Reproduce.} In \texttt{lean4/}, run \texttt{lake exe cache get},
\texttt{lake build}, then \texttt{lake env lean Check.lean}. Versions:
Lean/Mathlib \texttt{v4.33.0}. Independent finite checks are in
\texttt{reproduce.py}; they are not a substitute for the complete proof.

\textbf{Source.} \url{https://github.com/The-Last-Math-Competition/The-Last-Math-Competition/blob/28a22efac52c4a7abd4bba8a143b1c1a06b7e177/conjectures/00000000058.md}
\end{document}
